\documentclass[10pt,a4paper]{amsart}
\usepackage[margin=19mm]{geometry}
\usepackage{amsmath,amssymb,mathtools}
\usepackage[scr=rsfs,cal=euler]{mathalfa}
\usepackage{xfrac}
\usepackage[unicode,hidelinks,bookmarks=false]{hyperref}
\newcommand{\ie}{i.e.\ }
\newcommand{\cf}{cf.\ }
\newcommand{\resp}{respectively\ }
\newcommand{\Subsection}{\subsection}
\providecommand{\nobreakdash}{\nobreak\hbox{-}}
\setlength{\emergencystretch}{2em}
\multlinegap0pt
\usepackage{bm}
\usepackage{bbm}
\usepackage{dsfont}
\usepackage{xspace}
\usepackage{cancel}
\usepackage{mathtools}
\usepackage{amsthm}
\makeatletter
\@ifundefined{lemma}{\newtheorem{lemma}{Lemma}}{}
\@ifundefined{theorem}{\newtheorem{theorem}{Theorem}}{}
\makeatother
\def\RRR{\mathfrak{R}}
\def\SSS{\mathfrak{S}}
\def\ggg{\mathfrak{g}}
\title[Uniqueness of addition in Lie algebras revisited]{Uniqueness of addition in Lie algebras revisited}
\author{Ivan Arzhantsev}
\date{Source: arXiv:2401.06241v2 (CC BY 4.0)}
\begin{document}
\maketitle

\noindent\emph{Source and licence.} Uniqueness of addition in Lie algebras revisited. Version arXiv:2401.06241v2 (CC BY 4.0), distributed under the Creative Commons Attribution 4.0 International licence, as recorded in the arXiv licence metadata for this exact version and shown on the abstract page https://arxiv.org/abs/2401.06241v2. Full rights notice: licenses/arxiv-2401.06241-CC-BY-4.0.txt.\par\smallskip

\noindent\textit{Editorial scope.} Counted unit: the Theorem whose source label is \texttt{TCC} in the article (source0.tex lines 410--413, source label \texttt{TCC}), delivered together with the article's own complete original proof of it (source0.tex lines 415--472, one \texttt{proof} environment of 58 lines). No mathematics is written, repaired or re-derived: statement and proof are the article's own text line for line. Required context retained uncounted: TL (lines 383--387). Every definition, display and label of those lines is retained unchanged. Omitted from this package and not redistributed: 1--382, 388--409, 414--414, 473--892. Nothing inside a retained excerpt is omitted or altered: every retained body is delivered exactly as the article's lines are written, so no omission receipt of any kind accompanies this package. Citations: the retained excerpts carry no citation command at all, so no bibliography entry of the article is transcribed and no reference is redistributed. Reference adaptation: none was needed and none was made; every reference inside the delivered unit resolves inside it, so no unresolved reference or citation is produced. Printed slips kept literal: no printed word, symbol, hypothesis or number of the counted statement or of its proof was altered. Layout and class: the article's own class is \texttt{amsart} with packages including inputenc, fontenc, babel, amsmath, amssymb, amscd, amsthm, euscript, xy; the wrapper is the lane's pinned \texttt{amsart} editorial wrapper at 10pt on A4, which restates the theorem-like environments the retained text needs and adds the article's own macro definitions that the retained text needs (\texttt{\textbackslash RRR}, \texttt{\textbackslash SSS}, \texttt{\textbackslash ggg}), with extra wrapper packages (bm, bbm, dsfont, xspace, cancel, mathtools). The delivered numbering is the unit's own. Assets: no retained span contains a figure environment, a graphics-inclusion command, a PDF-page inclusion, a table or a TikZ picture; no image file is copied into this package, the package declares no asset dependency and no image file is redistributed with it. Licence: version arXiv:2401.06241v2 of this article is distributed under the Creative Commons Attribution 4.0 International licence, as recorded in the arXiv licence metadata for this exact version and shown on the abstract page https://arxiv.org/abs/2401.06241v2. A full rights notice is delivered at licenses/arxiv-2401.06241-CC-BY-4.0.txt. Characters: every retained excerpt is pure ASCII, so the delivered engine, XeLaTeX with its default Latin Modern text font, reports no missing character.
\begin{lemma} 
\label{TL}
Let $\alpha\colon\RRR\to\SSS$ be a commutator-preserving bijection of Lie rings. If  
for some $a,b\in\RRR$ we have $C(a)\cap C(b)=\{0\}$, then $\alpha(a+b)=\alpha(a)\oplus\alpha(b)$. 
\end{lemma}

\begin{theorem}
\label{TCC}
Any finite-dimensional Lie algebra over an infinite field satisfying the C-condition is a UA-Lie ring.
\end{theorem} 

\begin{proof} 
For a Lie algebra $\mathfrak{g}$ with C-condition we let
$$
C(\mathfrak{g})=\{a\in\mathfrak{g}\ | \  \exists \ b\in\mathfrak{g} \ \text{with} \ C(a)\cap C(b)=\{0\}\}.
$$

\begin{lemma}
For any $a\in C(\mathfrak{g})$ the set $U(a)=\{b\in\mathfrak{g} \mid C(a)\cap C(b)=\{0\} \}$ is open and dense in Zariski topology on $\mathfrak{g}$.
\end{lemma}
%
\begin{proof}
Consider the linear map
$$
\mathfrak{g}\to\mathfrak{g}\oplus\mathfrak{g}, \ \ \ \  x\mapsto ([x,a], [x,y]). 
$$
At the point $y=b$ the rank of this map is maximal. So the rank is maximal on a dense open subset $U(a)$ in $\mathfrak{g}$.
\end{proof}

Since the set $C(\ggg)$ is the union of subsets $U(a), a\in C(\ggg)$, the set $C(\ggg)$ is open and dense in $\ggg$ as well. 

\smallskip

Now let $\alpha\colon\ggg\to\SSS$ be a commutator-preserving bijection of Lie rings.

\bigskip

{\it Case 1.}\ Suppose $a, b\in C(\mathfrak{g})$. The intersection $U(b)\cap(U(a)-b)$ of Zariski open subsets is nonempty and we take an element $c$ from this intersection.
By Lemma~\ref{TL}, we obtain
$$
\alpha(a)\oplus\alpha(b)\oplus\alpha(c)=\alpha(a)\oplus\alpha(b+c)=\alpha(a+b+c)=\alpha(a+b)\oplus\alpha(c)\oplus\alpha(d),
$$
with some $d\in\ggg$, and $\alpha(a)\oplus\alpha(b)=\alpha(a+b)\oplus\alpha(d)$. As we have seen in the proof of Lemma~\ref{TL},
the equality $\alpha(a+b+c)=\alpha(a+b)\oplus\alpha(c)\oplus\alpha(d)$ implies that the elements $c$ and $d$ commute. 
Since $c$ runs through a dense open subset, we conclude that $d\in Z(\ggg)$. But the algebra $\ggg$
satisfies the C-condition, so its center is trivial and $d=0$. 

\bigskip

{\it Case 2.}\ Suppose $b\in C(\mathfrak{g})$. Taking an element $a_1$ from $C(\mathfrak{g})\cap(a-C(\mathfrak{g}))\cap(a+b-C(\mathfrak{g}))$
and letting $a_2=a-a_1$, we obtain $a_1, a_2\in C(\mathfrak{g})$ with $a_1+a_2=a$ and $a_2+b\in C(\mathfrak{g})$. Using Case~1, we get
$$
\alpha(a)\oplus\alpha(b)=\alpha(a_1)\oplus\alpha(a_2)\oplus\alpha(b)=\alpha(a_1)\oplus\alpha(a_2+b)=\alpha(a_1+a_2+b)=\alpha(a+b).
$$

\bigskip

{\it Case 3.}\  Let $a,b$ be arbitrary elements in $\ggg$. Taking $b_1$ from $C(\ggg)\cap(b-C(\ggg))$ and letting $b_2=b-b_1$, we obtain
$b=b_1+b_2$ with $b_1,b_2\in C(\mathfrak{g})$. Using Case~2, we get
$$
\alpha(a+b)=\alpha(a+b_1+b_2)=\alpha(a+b_1)\oplus\alpha(b_2)=
$$
$$
=\alpha(a)\oplus\alpha(b_1)\oplus\alpha(b_2)=\alpha(a)\oplus\alpha(b_1+b_2)=\alpha(a)\oplus\alpha(b).
$$
\smallskip

This completes the proof of Theorem~\ref{TCC}.
\end{proof} 

\end{document}
